% Part: methods
% Chapter: proofs
% Section: using-definitions

\documentclass[../../../include/open-logic-section]{subfiles}

\begin{document}

\olfileid[id]{mod}{prf}{def}

\olsection{Menggunakan Definisi}

Kita telah menyebutkan bahwa Anda harus memahami semua definisi yang
mungkin digunakan dalam bukti dan mampu menerapkannya dengan benar.
Hal ini sangat penting dan layak kita telaah sedikit lebih rinci.
Definisi digunakan untuk menyingkat sifat dan relasi sehingga kita
dapat membicarakannya dengan lebih ringkas. Singkatan yang
diperkenalkan disebut \emph{definiendum}, sedangkan apa yang
disingkatnya disebut \emph{definiens}. Dalam bukti, kita sering harus
kembali pada cara definiendum diperkenalkan, karena kita perlu
memanfaatkan struktur logis definiens (versi panjang yang disingkat
oleh istilah terdefinisi) agar dapat menuntaskan bukti. Dengan
menguraikan definisi, Anda memastikan bahwa Anda sampai ke inti tempat
kerja logis berlangsung.

Kita mulai dengan sebuah contoh. Andaikan Anda ingin membuktikan hal
berikut:

\begin{prop}
Untuk sebarang himpunan $A$ dan $B$, $A \cup B = B \cup A$.
\end{prop}

Agar dapat memulai bukti saja, kita perlu mengetahui apa artinya dua
himpunan identik; yaitu, kita perlu mengetahui apa arti ``$=$'' dalam
persamaan tersebut bagi himpunan. Dua himpunan didefinisikan identik
jika keduanya memiliki !!{element} yang sama. Jadi, definisi yang perlu
kita uraikan ialah:

\begin{defn}
Himpunan $A$ dan $B$ \emph{identik}, $A = B$, jika dan hanya jika setiap
!!{element} dari~$A$ merupakan !!a{element} dari~$B$, dan sebaliknya.
\end{defn}

Definisi ini menggunakan $A$ dan~$B$ sebagai pengganti bagi himpunan
sembarang. Yang didefinisikannya---\emph{definiendum}---adalah ungkapan
``$A = B$'', dengan memberikan syarat yang membuat $A = B$ benar.
Syarat ini---``setiap !!{element} dari~$A$ merupakan !!a{element}
dari~$B$, dan sebaliknya''---adalah \emph{definiens}.
\footnote{Dalam kasus khusus ini---dan sangat membingungkan!---ketika
  $A = B$, himpunan $A$ dan $B$ sebenarnya merupakan satu himpunan
  yang sama, walaupun kita menggunakan huruf yang berbeda untuknya di
  ruas kiri dan kanan. Namun, cara himpunan tersebut ditentukan dapat
  berbeda, dan itulah yang membuat definisi ini tidak sepele.}
Definisi tersebut menetapkan bahwa $A = B$ benar jika, dan hanya jika
(ungkapan ini dalam bahasa sumber disingkat menjadi ``iff''), syarat
tersebut terpenuhi.

Ketika menerapkan definisi, Anda harus mencocokkan $A$ dan $B$ dalam
definisi dengan kasus yang sedang dihadapi. Dalam kasus kita, agar $A
\cup B = B \cup A$ benar, setiap $z \in A \cup B$ juga harus berada
dalam $B \cup A$, dan sebaliknya. Ungkapan $A \cup B$ dalam proposisi
memainkan peran~$A$ dalam definisi, sedangkan $B \cup A$ memainkan
peran~$B$. Karena $A$ dan $B$ digunakan baik dalam definisi maupun
dalam pernyataan proposisi yang sedang kita buktikan, tetapi dengan
penggunaan yang berbeda, Anda harus berhati-hati agar keduanya tidak
tertukar. Misalnya, keliru jika mengira bahwa proposisi dapat
dibuktikan dengan menunjukkan bahwa setiap !!{element} dari~$A$
merupakan !!a{element} dari~$B$, dan sebaliknya---itu akan menunjukkan
bahwa $A = B$, bukan bahwa $A \cup B = B \cup A$. (Selain itu, karena
$A$ dan $B$ dapat berupa dua himpunan sebarang, Anda tidak akan
melangkah jauh: jika tidak ada apa pun yang diasumsikan mengenai $A$
dan~$B$, keduanya sangat mungkin merupakan himpunan yang berbeda.)

Di dalam bukti kita berurusan dengan gagasan-gagasan teori himpunan
seperti gabungan, sehingga kita juga harus mengetahui arti simbol
$\cup$ agar memahami bagaimana bukti harus dilanjutkan. Kadang-kadang,
penguraian suatu definisi melahirkan definisi lain yang juga perlu
diuraikan. Misalnya, $A \cup B$ didefinisikan sebagai
$\Setabs{z}{z \in A \text{ atau } z \in B}$. Jadi, jika Anda ingin
membuktikan bahwa $x \in A \cup B$, penguraian definisi $\cup$
memberi tahu Anda bahwa Anda harus membuktikan $x \in
\Setabs{z}{z \in A \text{ atau } z \in B}$. Sekarang Anda juga harus
mengingat bahwa $x \in \Setabs{z}{\dots z\dots}$ jika dan hanya jika
$\dots x\dots$. Dengan demikian, setelah notasi
$\Setabs{z}{\dots z \dots}$ diuraikan lebih lanjut, yang harus Anda
tunjukkan ialah: $x \in A$ atau $x \in B$. Jadi, ``setiap !!{element}
dari $A \cup B$ juga merupakan !!a{element} dari $B \cup A$''
sebenarnya berarti: ``untuk setiap $x$, jika $x \in A$ atau $x \in B$,
maka $x \in B$ atau $x \in A$.'' Jika semua definisi dalam proposisi
diuraikan sepenuhnya, kita melihat bahwa yang harus ditunjukkan adalah
hal berikut:

\begin{prop}
Untuk sebarang himpunan $A$ dan $B$: (a) untuk setiap $x$, jika $x \in
A$ atau $x \in B$, maka $x \in B$ atau $x \in A$, dan (b) untuk setiap
$x$, jika $x \in B$ atau $x \in A$, maka $x \in A$ atau $x \in B$.
\end{prop}

Yang penting ialah bahwa menguraikan definisi merupakan bagian yang
perlu dalam menyusun bukti. Melakukannya dengan benar kadang-kadang
sulit: Anda harus berhati-hati membedakan dan mencocokkan variabel
dalam definisi dengan suku dalam klaim yang sedang dibuktikan. Agar
berhasil, Anda harus mengetahui apa yang diminta oleh soal dan apa
arti semua istilah yang digunakan dalam soal---sering kali Anda perlu
menguraikan lebih dari satu definisi. Dalam bukti sederhana seperti
yang terdapat di bawah ini, penyelesaiannya hampir langsung mengikuti
definisi-definisi itu sendiri. Tentu saja, tidak selalu sesederhana
ini.

\begin{prob}
Andaikan Anda diminta membuktikan bahwa $A \cap B \neq \emptyset$.
Uraikan semua definisi yang muncul di sini, yaitu nyatakan kembali
pernyataan ini dengan cara yang tidak menyebut ``$\cap$'', ``='', atau
``$\emptyset$''.
\end{prob}

\end{document}
